\chapter*{Conclusion \& Perspectives}
\addcontentsline{toc}{chapter}{Conclusion \& Perspectives}

\section*{What was built, and what it demonstrates}

Over six sprints this project delivered \textbf{Tawer Management}, a two-application platform that consolidates
an agency's operational work into one authenticated, bilingual system: a stateless NestJS~11 REST API of
147~endpoints across 20~controllers, and a Next.js~16 / React~19 client, backed by PostgreSQL + pgvector. The four objectives set out in Section~1.6 were met. The breadth objective is visible in
the ten operational domains that Chapters~3 and~4 present sprint by sprint: projects and membership,
the agile backlog, a data-driven kanban, personal to-dos, attendance, calendar, reminders, multi-channel
notifications, infrastructure monitoring, and the AI copilot. The access-control objective is met by a
single centralized RBAC model that lets roughly 31 role types share one API across the two business units,
refined per-resource in the services rather than only at the guard (Chapter~3, Section~3.1). The maintainability
objective shows in the uniform four-layer, per-feature convention (controller → service → repository →
DTO) that gives every module the same structure, easy to read once the pattern is learned (Section~2.4). The
fourth objective --- the intelligent layer --- is met by the RAG subsystem of Chapter~4, and it is the one the
project achieves most fully.

Three things make the AI sprint the main contribution rather than a simple add-on. Its retrieval is
\textbf{permission-scoped in SQL}, so the same \texttt{projectId = ANY(\$allowedIds)} filter that governs the rest of
the platform also governs what the copilot can see. A CEO and an intern asking the identical question
retrieve from disjoint candidate sets, and the evaluation measured \textbf{zero cross-role leakage} (Sections~4.6
and~4.8). Its freshness is maintained by a \textbf{transactional-outbox boundary} processed by a locked sweeper, with a
self-healing nightly reconcile, so the domain never depends on the AI subsystem to save its data (Section~4.5).
And its quality is \textbf{measured rather than asserted}: an offline evaluation harness with committed gold
sets showed hybrid retrieval lifting MRR from 0.57 to 1.00 and Recall@1 from 0.30 to 1.00 on the
identifier gold set, while tying the already-saturated semantic set. That is the evidence that justifies the
hybrid lexical+vector design (Section~4.8). In a report built entirely from evidence, this is the one
feature backed by numbers.

\section*{Limitations \& Perspectives}

The platform is a functional system delivered under a strict six-sprint timeline: breadth was the priority,
and the hardening path ahead is short and well defined. The report named each gap in the sprint where
it appeared rather than saving it for the end, and \textbf{Annex~A} gathers every item into one ordered backlog,
each with its root cause and a planned fix. Together, the verified
issues fall into five themes:

\begin{itemize}
  \item \textbf{Security.} Before any non-local deployment, the platform needs the auth hardening named in Sprint~1 and
  the access-control fixes from Sprints~4 and~5: shorter token TTLs with an enforced \texttt{type} claim, request
  throttling and a global guard, \texttt{ValidationPipe(\{ whitelist: true \})} against mass-assignment, tighter
  CORS, moving frontend tokens off \texttt{localStorage} to HttpOnly cookies with edge auth re-enabled, and the
  \texttt{deleteUserByAdmin}/IDOR fixes (Annex~A; Section~3.1.6).
  \item \textbf{Correctness.} A small set of concrete, code-verified defects, each with its root cause and a fix in
  Annex~A: the check-in dead zone, server-side archived filtering, the racy \texttt{generateTaskKey},
  reminder recurrence, \texttt{checkHttp} on a schemeless domain, event ownership on edit, and the AI
  comment-delete freshness gap.
  \item \textbf{Maintainability.} The duplicated executive-RBAC helpers and the copy-pasted notification fan-out
  consolidate behind a shared \texttt{ProjectAccessService} and a single notification dispatcher; a set of unused
  structures (the single-value \texttt{Language} enum, the dead \texttt{RedisModule}/\texttt{RolesGuard}/\texttt{UserManager}/
  \texttt{Task.statusType}, stale branding) should be connected or removed rather than left misleading.
  \item \textbf{Testing.} Backend behavioural coverage is near-zero outside the projects module (the 854-line
  \texttt{projects.service.spec.ts}, cited as genuine coverage in Section~3.2.5) and the frontend property-based suites;
  the highest-value missing tests are the ones the bugs above name, and a business-day-boundary test would
  have caught the check-in dead zone.
  \item \textbf{DevOps.} No CI/CD, app Dockerfiles not yet orchestrated by compose, uploads stored on a
  non-persistent local filesystem, and a frontend \texttt{BACKEND\_ADDRESS} hard-coded to \texttt{localhost}, so the
  platform runs locally but is not yet ready to deploy elsewhere.
\end{itemize}

None of these weaken the architecture. They are the finishing work that separates a working platform
from a production system. Naming them precisely, with fixes gathered in Annex~A, is what makes
the prioritized path below concrete.

\section*{Perspectives: prioritized future work}

The roadmap starts with the most damaging issues:

\begin{enumerate}
  \item \textbf{Security first.} Shorten token TTLs and enforce the \texttt{type} claim, add \texttt{@nestjs/throttler} plus a
   global auth guard, enable \texttt{whitelist} validation, restrict CORS, move the frontend to HttpOnly cookies
   with edge auth re-enabled, and fix \texttt{deleteUserByAdmin} and the two IDOR writes. These are the changes
   that gate any real deployment.
  \item \textbf{Correctness.} Key WorkDays on an explicit \texttt{businessDate} column to fix the dead zone; replace
   \texttt{generateTaskKey} with a monotonic per-project counter; enqueue comment edits and deletes into the AI
   outbox so the copilot never cites content that is gone; and scope \texttt{ProjectContent}/\texttt{SprintContent}
   uniqueness by project.
  \item \textbf{Quality \& maintainability.} Extract a shared \texttt{ProjectAccessService}, split the large \texttt{TasksService},
   introduce a central notification dispatcher to replace the duplicated four-channel fan-out, and add an
   axios interceptor to remove the \textasciitilde60× duplicated frontend auth/refresh logic. Add the cheap unit tests
   around the AI subsystem's pure functions (RRF fusion, weighted percentiles, citation parsing).
  \item \textbf{DevOps.} Add CI (lint + build + \texttt{prisma validate} + the existing suites), a multi-stage Dockerfile,
   persistent upload storage, and a runtime-configurable \texttt{BACKEND\_ADDRESS}.
  \item \textbf{AI depth.} Make the confidence gate hybrid-aware so the answer path inherits the proven keyword improvement,
   add a retention/redaction policy for the plaintext \texttt{CopilotQueryLog}, and refresh the answer-quality
   evaluation figures.
  \item \textbf{i18n.} Either populate the \texttt{Language} enum and connect the content tables end to end, or remove
   the unused split, since the current half-way state is the misleading option.
\end{enumerate}

\section*{Closing}

Tawer Management set out to consolidate fragmented team-management workflows into one coherent, localized
application, and to add an evidence-grounded AI layer on top of the organization's own data. It meets that
goal: ten operational domains sit behind one RBAC-governed API, and the RAG copilot answers only from
content the user is allowed to see, with citations and a measured quality score. The main engineering
strengths are real. The architecture is uniformly layered, access control is centralized, and the RAG
pipeline is designed the way a production system would be. The remaining work is hardening, and it is prioritized in Annex~A
rather than left vague. The retrieval numbers are the clearest evidence that the hardest part already
works: on the identifier gold set, MRR rose from 0.57 to 1.00 and Recall@1 from 0.30 to 1.00, with no
cross-role leakage.
